\scanpage{82}
% archive foliation: 81
$Z$ constructible dans $X$ qu. compact,
\inserted{\textcircled{\unsure{ou}}}
\unsure{tt pt} $\in$ une \unsure{petite fam. de voisinages rétrocompacts dans} $X$.
\deleted{$Z$ loc. constructible}
\[
\left\{\begin{array}{l}
Z\ \text{qu. cpt et}\\
\forall x\in Z,\ \exists\ \text{un voisinage }U_x\text{ de }x\text{ dans }X,
\text{ tel que }Z\cap U_x\text{ constr. dans }U_x
\end{array}\right.
\]
(alors $\unsure{U}\longrightarrow U'\subset U$, $U'\ni x$, a la $=$ propriété)

nécessité : évident [car $Z$ rétrocompact dans $X$, \deleted{\unsure{dans}}
\unsure{et} on peut prendre $U=X$.

Suffisance. $Z$ est recouvert par un nb fini des $U_x$, soient $U_i$.
Les $U_i\cap Z$ sont constructibles dans $U_i$, donc dans $X$,
donc $Z=\bigcup Z_i$ est constr.

\uline{Corollaire} Supposons $Z$ loc. fermé.
Alors $Z$ \deleted{\unsure{loc. constructible}}
ssi $Z$ est \uline{quasi-compact}, et tt $x\in Z$ a un voisinage ouvert
$U$ dans $X$ tel que $Z\cap U$ soit \uline{fermé} et
\uline{constructible} dans $U$, i.e. $U-Z\cap U$ rétrocompact dans $U$
[i.e. $U-Z\cap U$ quasi-compact,
\inserted{(pour $U$ qu-compact, et si on)}
\unsure{suppose que dans} $X$ l'intersection de deux ouverts quasi-compacts
est qu. cpt)]

\uline{Exemple} Si $X$ est un préschéma quasi-séparé quasi-compact
(p. ex. un schéma quasi-compact, ou un schéma affine) \deleted{\unsure{donc}}
$Z$ loc. fermé dans $X$, alors $Z$ constructible $\Longleftrightarrow$
$Z$ qu. cpt, et tt $x\in Z$ a un voisinage ouvert affine $U$ tel que
$Z\cap U$ \inserted{soit défini par un} \uline{syst. fini} \uline{d'équations}.

% The reversed typewriting showing through this sheet is on the separately supplied reverse, scan 83, where it is transcribed.
